% --- 1. INFORMAZIONI GENERALI ---
\section{Informazioni Generali}
\begin{itemize}
    \item \textbf{Luogo/Piattaforma:} Server Discord\textsubscript{G}
    \item \textbf{Data:} 2026-06-23
    \item \textbf{Orario di inizio:} 14:30
    \item \textbf{Orario di fine:} 15:30
    \item \textbf{Responsabile\textsubscript{G}:} Armando Moda Scarati
    \item \textbf{Scriba:} Sofia De Blasi
\end{itemize}

\vspace{0.5cm}
\textbf{Partecipanti presenti:}
\begin{table}[ht]
    \centering
    \renewcommand{\arraystretch}{1.2}
    \begin{tabularx}{\textwidth}{X l}
        \toprule
        \textbf{Nome e Cognome} & \textbf{Matricola} \\
        \midrule
        Sofia De Blasi & 2111014 \\
        Roberto Mariano Doroftei & 2111031 \\
        Filippo Idiometri & 2035617 \\
        Francesco Marcon & 2110990 \\
        Armando Moda Scarati & 2082864 \\
        Anton Ricardo Rupi & 2054304 \\
        \bottomrule
    \end{tabularx}
\end{table}

\vspace{0.5cm}
\textbf{Partecipanti assenti:}
\begin{table}[ht]
    \centering
    \renewcommand{\arraystretch}{1.2}
    \begin{tabularx}{\textwidth}{X l}
        \toprule
        \textbf{Nome e Cognome} & \textbf{Matricola} \\
        \midrule
        Nessuno & - \\
        \bottomrule
    \end{tabularx}
\end{table}

% --- 2. ORDINE DEL GIORNO ---
\section{Ordine del Giorno}

Gli argomenti previsti per la riunione odierna sono:

\begin{enumerate}
    \item Sprint\textsubscript{G} review e retrospective,
    \item Aggiornamento dei documenti di progetto (PdP, PdQ) e delle slide di presentazione,
    \item Discussione sulle criticità infrastrutturali e comunicazioni con la proponente,
    \item Analisi preliminare per la Product Baseline\textsubscript{G} (architettura, CI/CD\textsubscript{G}, metriche),
    \item Pianificazione dello sprint successivo e allocazione dei ruoli.
\end{enumerate}


% --- 3. RESOCONTO E DISCUSSIONE ---
\section{Resoconto della Riunione}

\subsection{Sprint review\textsubscript{G} e retrospective}
Il gruppo ha effettuato un rapido confronto sull'andamento dello sprint precedente. Non sono state riscontrate problematiche interne rilevanti e i documenti risultano in regola. Tuttavia, a livello organizzativo, è stato rilevato un rallentamento generale dovuto alla sovrapposizione con i vari impegni accademici dei membri. Dal punto di vista tecnico, è emersa una criticità bloccante legata alla latenza di risposta dei modelli LLM\textsubscript{G} resi disponibili dall'infrastruttura dell'azienda proponente. A tal proposito, il Responsabile si incaricherà di redigere un'email formale diretta all'azienda chiedendo un incontro in cui il gruppo illustrerà il Proof of Concept\textsubscript{G} e chiederà delucidazioni o alternative in merito al problema infrastrutturale.

\subsection{Aggiornamento dei documenti e presentazioni}
Le attività di chiusura per la fase attuale procedono senza intoppi. Il ruolo\textsubscript{G} dell'Amministratore\textsubscript{G} si farà carico dell'aggiornamento del Piano di Progetto\textsubscript{G}, inserendo il consuntivo\textsubscript{G} per lo Sprint 11, aggiornando le percentuali del preventivo a finire e allineando l'organigramma. Contestualmente, interverrà sul Piano di Qualifica\textsubscript{G} per perfezionare le descrizioni e i diagrammi. Infine, dovrà assicurarsi di riportare e aggiornare correttamente il monte ore all'interno delle slide di presentazione destinate al colloquio con il prof. Vardanega.

\subsection{Preparazione alla Product Baseline}
In attesa degli esiti ufficiali e del colloquio, il team ha concordato di non interrompere i lavori, sfruttando il periodo di transizione per iniziare a ragionare sulla futura architettura del prodotto. Tutti i membri avvieranno uno studio preliminare che servirà da base per la stesura della Specifica Tecnica. Nello specifico, la ricerca si concentrerà sull'architettura logica e di deployment, sull'individuazione dei design pattern, delle dipendenze e delle interfacce. Parallelamente, verrà analizzato il setup di un ambiente CI/CD (con potenziale uso di GitHub Actions\textsubscript{G}) per automatizzare le fasi di build e testing, e si decideranno in anticipo le metriche di code quality e i relativi strumenti di misurazione\textsubscript{G} da integrare.

% --- 4. DECISIONI PRESE ---
\section{Decisioni Prese}

\renewcommand{\arraystretch}{1.3}
\vspace{-0.3cm}
\noindent
\begin{longtable}{c p{12cm}}
    \toprule
    \textbf{Codice} & \textbf{Decisione} \\
    \midrule
    \endfirsthead

    \toprule
    \textbf{Codice} & \textbf{Decisione} \\
    \midrule
    \endhead

    \midrule
    \multicolumn{2}{r}{\textit{Continua nella pagina successiva}} \\
    \midrule
    \endfoot

    \bottomrule
    \endlastfoot

    \textbf{VI-19.1} & {Pianificata la comunicazione ufficiale alla proponente in merito alle problematiche di latenza dei modelli LLM e per la visione del POC\textsubscript{G}.} \\
	\textbf{VI-19.2} & {Avvio di una fase di studio collettiva per la predisposizione degli strumenti di CI/CD, architettura e metriche necessarie per la futura Specifica Tecnica.}
\end{longtable}

% --- 5. AZIONI (TODO) ---
\section{Azioni da Svolgere (To-Do)}

\renewcommand{\arraystretch}{1.3}
\begin{longtable}{p{2.3cm} c p{4.5cm} p{3cm} l}
    \toprule
    \textbf{Codice} & \textbf{Rif. Decisione} & \textbf{Descrizione Task\textsubscript{G}} & \textbf{Assegnatario} & \textbf{Scadenza} \\
    \midrule
    \endfirsthead

    \toprule
    \textbf{Codice} & \textbf{Rif. Decisione} & \textbf{Descrizione Task} & \textbf{Assegnatario} & \textbf{Scadenza} \\
    \midrule
    \endhead

    \midrule
    \multicolumn{5}{r}{\textit{Continua nella pagina successiva}} \\
    \midrule
    \endfoot

    \bottomrule
    \endlastfoot

    \ghissue{253} & VI-16.1 & Stesura e invio email a Zucchetti\textsubscript{G} per POC e criticità latenza modelli & Armando M. Scarati & 2026-06-30 \\[4pt]
	
	\ghissue{255} & - & Aggiornamento consuntivo Sprint 11, preventivo e organigramma nel PdP & Roberto M. Doroftei & 2026-06-30 \\[4pt]
	
	\ghissue{256} & - & Aggiornamento diagrammi e relative descrizioni all'interno del PdQ & Roberto M. Doroftei & 2026-06-30 \\[4pt]
	
	\ghissue{257} & - & Integrazione e aggiornamento del riepilogo ore nelle slide per il prof. Tullio & Roberto M. Doroftei & 2026-06-30 \\[4pt]
	
	& VI-16.2 & Studio preliminare su architettura, automazioni CI/CD e metriche di qualità & Tutti i membri & 2026-06-30 \\[4pt]
	
	\ghissue{254} & - & Redazione del verbale interno della riunione odierna (VI N. 19) & Armando M. Scarati & 2026-06-30 \\[4pt]
\end{longtable}

% --- 6. PROSSIMA RIUNIONE ---
\section{Prossima Riunione}
\begin{itemize}
	\item \textbf{Data:} 2026-06-30
	\item \textbf{Ora:} 14:30
	\item \textbf{OdG preliminare\textsubscript{G}:}
	\begin{itemize}
		\item Verifica dei task assegnati per lo Sprint 11;
		\item Allineamento sulle comunicazioni ricevute dall'azienda proponente;
		\item Condivisione dei risultati dello studio preliminare sull'architettura e CI/CD;
		\item Pianificazione dello sprint successivo e assegnazione dei ruoli.
	\end{itemize}
\end{itemize}

